\documentclass[12pt,utf8]{beamer}


%Wichtige Standard Pakete!
\usepackage[utf8]{inputenc}
\usepackage[ngerman]{babel}
\usepackage{xcolor}
\usepackage{graphicx}
\usepackage{amsmath}
\usepackage{mathtools}
\usepackage{amssymb}
\usepackage{amsthm}
\usepackage{listings}
\usepackage{mathabx}
\usepackage{tikz,pgf}
\usepackage{tikz-cd}
\usepackage{pgf}
\usepackage{algorithm, algorithmic}
\renewcommand{\algorithmicrequire}{\textbf{Eingabe:}}
%\renewcommand{\algorithmicensure}{\textbf{Ausgabe:}}
\usetikzlibrary{matrix,arrows}
\usepackage[autostyle=true,german=quotes]{csquotes}
\usepackage[all]{xy}
\usepackage{stmaryrd}
\usepackage{framed}
\usepackage{bbm}
\usepackage{cite}
\usepackage{etoolbox}
\usepackage{hyperref}
\usepackage{enumitem}
\newcommand\mySlash[2]{\ensuremath{%
		\!\sideset{_#1}{\!^#2}{\mathop\backslash}}}
\apptocmd{\sloppy}{\hbadness 10000\relax}{}{}

%Nummerierung der Folien
\beamertemplatefootpagenumber 


%Programmiersprache festlegen
\definecolor{dkgreen}{rgb}{0,0.6,0}
\definecolor{gray}{rgb}{0.5,0.5,0.5}
\definecolor{mauve}{rgb}{0.58,0,0.82}

\lstdefinestyle{myScalastyle}{
  frame=tb,
  language=scala,
  aboveskip=3mm,
  belowskip=3mm,
  showstringspaces=false,
  columns=flexible,
  basicstyle={\small\ttfamily},
  numbers=left,
  numberstyle=\tiny\color{gray},
  keywordstyle=\color{blue},
  commentstyle=\color{dkgreen},
  stringstyle=\color{mauve},
  frame=single,
  breaklines=true,
  breakatwhitespace=true,
  tabsize=2,
}


% Definiere div-Operator äquivalent zu mod
\makeatletter
\newcommand*{\bdiv}{%
  \nonscript\mskip-\medmuskip\mkern5mu%
  \mathbin{\operator@font div}\penalty900\mkern5mu%
  \nonscript\mskip-\medmuskip
}
\newcommand{\testleftlong}{\longleftarrow\!\shortmid}

\makeatother
\usepackage{rotating}
\newcommand\tabrotate[1]{\begin{turn}{90}\rlap{#1}\end{turn}}
\newcommand{\invertediota}{\begin{sideways}%
		\begin{sideways}$\iota$\end{sideways}\end{sideways}}
% Benutze theoremstyles von Thomas, ein klein wenig verändert
\newtheoremstyle{break-italic}% name
  {1ex}%      Space above, empty = `usual value'
  {}%      Space below
  {\itshape}% Body font
  {0cm}%         Indent amount (empty = no indent, \parindent = para indent)
  {\bfseries}% Thm head font
  {}%        Punctuation after thm head
  {\newline }% Space after thm head: \newline = linebreak
  {}%         Thm head 
\newtheoremstyle{break-roman}% name
  {1ex}%      Space above, empty = `usual value'
  {}%      Space below
  {\normalfont}% Body font
  {}%         Indent amount (empty = no indent, \parindent = para indent)
  {\bfseries}% Thm head font
  {}%        Punctuation after thm head
  {\newline }% Space after thm head: \newline = linebreak
  {}%         Thm head spec
  
\theoremstyle{break-italic} 
\newtheorem{defe}[theorem]{Definition}

\newcommand{\einschraenkung}{\,\rule[-5pt]{0.4pt}{12pt}\,{}}




%Für den Header notwendig!
\usepackage[percent]{overpic}

%Sourcecode
\usepackage{listings}

%Einbinden des Themes
%\input{Theme/ubuntuuserstheme.tex}

%Mathe
\usepackage{amsthm}

%Standard Angaben
\author{Jan Albert}
\title{Funktionale Programmierung in Scala}
\date{\today}
%\titlegraphic{\includegraphics[scale=0.04]{book2.jpg}}

\begin{document}


	%Titelseite
	\begin{frame}
		\titlepage
	
	\end{frame}
	
	%Inhaltsverzeichnis
	\begin{frame}
		\frametitle{Inhaltsverzeichnis}
		\tableofcontents
	\end{frame}
	
	%Buch
		\section[Buch]{Buch}
		\begin{frame}[fragile]
		\frametitle{Buch}
	    
\begin{columns}[c]
  \begin{column}{0.55\textwidth}
    Functional Programming in Scala\\
    \leavevmode \\
    Paul Chiusano,
    Runar Bjarnason\\
    \leavevmode \\
    Manning, 2014
  \end{column}
  \begin{column}[T]{0.55\textwidth}
    \includegraphics[scale=0.06]{book2.jpg}
  \end{column}
\end{columns}
		
		
\end{frame}
	
	%Einfuehrung
		%Einfuehrung
	\section[Section]{Was ist Funktionale Programmierung?}
		\begin{frame}
			\frametitle{Was ist Funktionale Programmierung?}
				Idee:
				Verwendung "reiner Funktionen"  (keine Seiteneffekte) \\
				\leavevmode \\
				Beispiele für Seiteneffekte:
				\begin{itemize}
				\item[•] verändern/modifizieren einer Variable
				\item[•] verändern/modifizieren einer Datenstruktur
				\item[•] eine Exception werfen
				%\item[•] Konsoleneingabe/-ausgabe
				%\item[•] lesen/schreiben aus/von einer Datei
				\end{itemize}
		\end{frame}
	
	%Definitionen
		\section[Section]{Definitionen}
	\begin{frame}
		\frametitle{Reine Funktion}
		\begin{definition}
			Eine \emph{reine Funktion} mit Eingabetyp $A$ und Ausgabetyp $B$ (Schreibweise: $A				\Rightarrow B$)
			ist eine Berechnung, welche jeden Wert $a$ vom Typ $A$ genau einen Wert 			$b$ vom Typ $B$ zuordnet, sodass $b$ nur aus dem Wert von $a$ bestimmt 				wird.
		\end{definition} 
		\leavevmode \\
		Beispiele: 
		\begin{itemize}
		\item[•] Funktion intToString 
		\item[•] Addition ganzer Zahlen
		\end{itemize}
	\end{frame}
	

	%Ausdruck
		\section*{}
		\begin{frame}
			\frametitle{Ausdruck}
			\begin{definition}
				Jeder Teil eines Programms, welcher zu einem Ergebnis 								zusammengefasst werden kann, d. h. alles was man in den 
				Scala-Interpreter tippt und ein Ergebnis liefert,
				nennt man einen \emph{Ausdruck}.
		\end{definition} 
		\leavevmode \\
		Beispiel:  $2+3$	
		\end{frame}
	
	%RT
		\section*{}
		\begin{frame}
			\frametitle{Referenziell Transparent}
			\begin{definition}
				Ein Ausdruck $e$ ist \emph{Referenziell Transparent (RT)},
				wenn für alle Programme $p$, alle Vorkommnisse von
				$e$ in $p$ durch das Ergebnis von $e$ ersetzt werden
				können, ohne die Bedeutung von $p$ zu ändern.
				Eine Funktion ist rein, wenn der Ausdruck $f(x)$ 			  
				referenziell transparent für alle referenziell
				transparenten $x$ ist.
			\end{definition} 
		\end{frame}
	
	%Beispiel1
		\section[Section]{Beispiel mit Seiteneffekt}
		\begin{frame}[fragile]
		\frametitle{Beispiel mit Seiteneffekt}
		\begin{lstlisting}[style=myScalastyle]
class Cafe {
	def buyCoffee(cc: CreditCard): Coffee = {
		val cup = new Coffee()
		cc.charge(cup.price)
		cup
	}
}
\end{lstlisting}
\end{frame}
		
			
	
	%Beispiel2
		\section*{}
		\begin{frame}[fragile]
		\frametitle{Beispiel ohne Seiteneffekt}
		\begin{lstlisting}[style=myScalastyle]
class Cafe {
	def buyCoffee(cc: CreditCard): (Coffee, Charge) = {
		val cup = new Coffee()
		(cup, Charge(cc, cup.price))
	}
}
\end{lstlisting}
\end{frame}
	
	%Spielerei
		\section[Section]{Spielerei}
		\begin{frame}[fragile]
		\frametitle{Buch}
	    
\begin{columns}[c]
  \begin{column}{0.55\textwidth}
    Functional Programming in Scala\\
    \leavevmode \\
    Paul Chiusano,
    Runar Bjarnason\\
    \leavevmode \\
    Manning, 2014
  \end{column}
  \begin{column}[T]{0.55\textwidth}
    \includegraphics[scale=0.06]{book2.jpg}
  \end{column}
\end{columns}
		
		
\end{frame}
		
			
	
	%Quellen
	\section[Section]{Quellen}
	\begin{frame}
  		\frametitle{Quellen}    
  		\begin{thebibliography}{10}    
  		\beamertemplatebookbibitems
  		\bibitem{Autor1990}
    		Paul Chiusano,
    		Runar Bjarnason
   		 \newblock {\em Functional Programming in Scala}
    	\newblock Manning, 2014.
  		\beamertemplatearticlebibitems
  		\bibitem{Jemand2000}
    	S.~Jemand.
    	\newblock On this and that.
    	\newblock {\em Journal of This and That}, 2(1):		50--100, 2000.
  		\end{thebibliography}
	\end{frame}
	
	%Danksagung
	\section[Section]{Danksagung}
	\begin{frame}
		\vfill
		\begin{center}\begin{Huge}Vielen Dank \\[10pt]
			für eure Aufmerksamkeit.\end{Huge}\vfill
		\end{center}
		\vfill
	\end{frame}
	
\end{document}